\section{Exploratory Results}
\label{sec:exploratory_results}

This section presents the first descriptive patterns in the evolution of housing development over time within the city.
The goal is to assess whether the raw data move in the direction the institutional story implies, and whether the pattern is concentrated on the housing margins where discretionary review should matter most.

\subsection{Raw Trends by Treatment Tercile}

Figure~\ref{fig:tercile_trends} groups community districts into low, middle, and high terciles of the within-borough homeownership treatment, then plots raw units built from 1970 through 2025.
The displayed series are centered 3-year moving averages of annual unit counts.
These raw counts are useful for the first descriptive question: where did new units actually appear?

The large-building panel is the main focus.
Before the reform, high-treatment districts produce a substantial number of large-building units.
After 2010, the ranking reverses: low-homeownership community districts produce many more large-building units than high-homeownership community districts.
The total-units panel shows a similar pattern.
The 1--4 unit margin, by contrast, does not exhibit the same late decline in the regression analogue below, which is reassuring because small buildings are less likely to be the margin where member deference operates.

\begin{figure}[htbp]
    \centering
    \includegraphics[page=1,width=0.86\textwidth]{input/cd_homeownership_long_units_raw_units_plots.pdf}
    \caption{Housing Production by Homeownership Tercile}
    \label{fig:tercile_trends}
\figurenotes{Figure plots centered 3-year moving averages of raw annual housing-production units by tercile of the within-borough standardized 1990 homeownership treatment. The dashed line marks 1990, the year in which homeownership exposure is measured. The unit is the community district. Construction outcomes cover 1970--2025 and are measured with the 25v4 MapPLUTO yearbuilt proxy.}
\end{figure}

\subsection{Brooklyn Case Study}

Brooklyn is a useful way to see the treatment variation within a single borough, without changing the citywide design.
Figure~\ref{fig:brooklyn_case_study} orders Brooklyn's 18 community districts by the same 1990 homeowner exposure and places later construction outcomes alongside the treatment.
This is descriptive, but it makes the empirical margin transparent: post-2020 large-building production is concentrated away from the homeowner-heavy community districts.

\begin{figure}[htbp]
    \centering
    \includegraphics[page=2,width=0.86\textwidth]{input/cd_homeownership_long_units_raw_units_plots.pdf}
    \caption{Brooklyn Homeownership Gradient by Housing Margin}
    \label{fig:brooklyn_case_study}
\figurenotes{Figure orders Brooklyn's 18 community districts by within-borough standardized 1990 homeownership and plots later construction margins alongside the treatment. Construction outcomes are raw 2020--2025 unit totals, measured with the 25v4 MapPLUTO yearbuilt proxy.}
\end{figure}

\subsection{Event-Study Results}
\label{subsec:regression_analogues}

I next estimate a community-district regression analogue of Figure~\ref{fig:tercile_trends}.
The preferred outcome is raw units in 5+ unit new buildings --- a margin broad enough to be well-powered, but one that still excludes the 1--4 unit range where discretionary review should matter less.
I plot the 1--4 unit series alongside as a placebo.

The event-study specification is
\begin{equation}
Y_{dbt}
= \alpha_d + \lambda_{b p(t)}
+ \sum_{p \neq p_0} \beta_p \left(T_d \times \mathbf{1}\{p(t)=p\}\right)
+ \sum_{p \neq p_0} C_d^{\prime}\gamma_p \mathbf{1}\{p(t)=p\}
+ \varepsilon_{dbt},
\label{eq:homeownership_event_study}
\end{equation}
where \(d\) indexes community districts, \(b\) boroughs, and \(p(t)\) period bins.
The omitted period \(p_0\) is 1985--1989.
The event-study figure uses five-year bins from 1970--1974 onward, with 2020--2025 as the final partial bin.
The treatment \(T_d\) is the 1990 community-district homeownership share, standardized within borough.
Controls include log 1990 occupied units, 1990 median household income, 1990 vacancy rate, and 1970--1988 raw production on the same outcome scale, interacted with the period bins.
Fixed effects absorb permanent community-district differences and borough-period housing cycles; standard errors are clustered by community district.

\begin{figure}[htbp]
    \centering
    \includegraphics[page=1,width=0.86\textwidth]{input/cd_homeownership_long_units_event_coefficients_raw_units_5yr_bins.pdf}
    \caption{Event-Study Gradients by Housing Margin}
    \label{fig:homeownership_event_study_5_plus}
\figurenotes{Figure plots coefficients on within-borough standardized 1990 homeownership from Equation~\ref{eq:homeownership_event_study}. Outcomes are raw annual units in 1--4 unit and 5+ unit new buildings. The unit is the community district. Construction outcomes are measured with the 25v4 MapPLUTO yearbuilt proxy. The omitted period is 1985--1989. The event-study bins are five-year periods from 1970--1974 through 2015--2019, plus a final 2020--2025 bin. Controls include log 1990 occupied units, 1990 median household income, 1990 vacancy rate, and 1970--1988 raw production on the same outcome scale, interacted with period bins. Standard errors are clustered by community district.}
\end{figure}

Figure~\ref{fig:homeownership_event_study_5_plus} does not show a clean break exactly at 1990.
For 5+ unit buildings, the homeowner gradient is negative in several earlier bins but becomes much larger after 2010, especially in the 2015--2019 and 2020--2025 periods.
The 1--4 unit placebo shows no comparable late decline.

\input{input/cd_homeownership_long_units_long_difference_raw_units.tex}

Table~\ref{tab:cd_homeownership_long_units_long_difference_raw_units} reports the corresponding long-difference estimates.
The results are consistent with the event study: the pre-1990 placebo comparison and the early post windows are imprecise, while the 2010s and especially the 2020--2025 long differences are negative.
These estimates do not capture the total citywide effect of Morris or member deference; they measure whether post-2010 production shifted away from homeowner-heavy community districts within the same borough.
